% !TeX program = xelatex
% UTF-8. Compile from this directory: xelatex shock.tex
% Source: Google Drive notebook/shock; photographs in chronological order.
% pic01.jpg = 20261016_044751075.jpg
% pic02.jpg = 20261016_044756378.jpg
% pic03.jpg = 20261016_044759457.jpg
% pic04.jpg = 20261016_044808551.jpg
% pic05.jpg = 20261016_044812030.jpg
% pic06.jpg = 20261016_044820535.jpg
% pic07.jpg = 20261016_044823766.jpg
% pic08.jpg = 20261016_044830534.jpg
% pic09.jpg = 20261016_044834601.jpg
% pic10.jpg = 20261016_044840615.jpg
% pic11.jpg = 20261016_044843860.jpg
% pic12.jpg = 20261016_044849594.jpg
% pic13.jpg = 20261016_044853197.jpg
% pic14.jpg = 20261016_044858556.jpg
\documentclass[12pt,a4paper]{article}
\usepackage[margin=22mm]{geometry}
\usepackage{amsmath,amssymb}
\usepackage{graphicx}
\usepackage{xepersian}
\IfFontExistsTF{Amiri}{\settextfont{Amiri}}{\settextfont{DejaVu Sans}}
\setlatintextfont{DejaVu Serif}
\setlength{\parindent}{0pt}
\setlength{\parskip}{5pt}
\setlength{\emergencystretch}{3em}
\newcommand{\ShE}[1]{\lr{#1}}
\newcommand{\ShM}[1]{\lr{\(#1\)}}
\newcommand{\ShU}[1]{\textbf{[خوانش نامطمئن: #1]}}
\newcommand{\ShPage}[1]{\clearpage\section*{تصویر #1 ـ شوک، ۸۵}}
\begin{document}
\begin{center}
{\LARGE\bfseries شوک}\par
\ShE{Nelson 22}\par
فصل ۸۵
\end{center}
متن به ترتیب تصاویر رونویسی شده است. فلش‌ها و شاخه‌های دست‌نویس به صورت رابطه یا فهرست آمده‌اند؛ خوانش‌های نامطمئن مشخص شده‌اند. تصاویر اصلی برای مقابله در پایان سند قرار دارند. سربرگ چاپی همهٔ برگه‌ها «موضوع» و «تاریخ» است؛ موضوع «شوک ـ ۸۵»، تاریخ خالی با دو ممیز، و شماره‌های پایین برگه‌ها ۱ تا ۱۴ هستند.

\section*{تصویر ۱}
شوک اگر در مراحل اولیه درمان نشود:
\begin{latin}
\begin{center}
Multiple Organ Dysfunction Syndrome (MODS)\\
\(\longrightarrow\) Cardiovascular Collapse \(\longrightarrow\) Death
\end{center}
\end{latin}
تظاهرات شوک ناشی از افت فشار و \ShE{mismatch} بین \ShE{supply} و \ShE{demand} است.

\subsection*{تظاهرات زودرس شوک}
\begin{itemize}
\item تاکی‌کاردی.
\item \ShE{MSC}: آژیتاسیون، لتارژی، تحریک‌پذیری.
\item تاکی‌پنه: عموماً ثانویه به اسیدوز.
\item زمان پرشدن مویرگی \ShE{(CFT)} غیرطبیعی.
\item از بین رفتن پالس‌های محیطی.
\item الیگوری.
\item مارکرهای متابولیک.
\item اسیدوز.
\item هیپرلاکتاتمی.
\item افت \ShM{SpO_2}.
\end{itemize}
زمانی که \ShM{BP\downarrow} رخ دهد، یعنی بیمار دچار فاز دکامپنسیتهٔ شوک است.

\ShE{MODS}: اختلال عملکرد در حداقل ۲ ارگان.

\ShPage{۲}
\subsection*{انواع شوک}
هایپوولمیک، کاردیوژنیک، آبستراکتیو، دیستریبیوتیو و دیسوسییتیو.

\subsection*{شوک هایپوولمیک}
شایع‌ترین شوک.

\subsection*{شوک کاردیوژنیک}
\begin{itemize}
\item در بیمار با بیماری قلبی مادرزادی (قبل و بعد جراحی).
\item در انواع کاردیومیوپاتی‌ها.
\item میوکاردیت حاد.
\item \ShE{MIS-C} (سندرم التهاب چندسیستمی در اطفال).
\end{itemize}
\subsection*{شوک انسدادی}
\begin{itemize}
\item وجود ضایعاتی که جلوی جریان خون از قلب را می‌گیرد.
\item تامپوناد، تنشن پنوموتوراکس، \ShE{PE}، \ShE{AS} شدید، کوآرکتاسیون آئورت.
\end{itemize}
\subsection*{شوک توزیعی}
\begin{itemize}
\item وازودیلاتاسیون غیرطبیعی؛ خون در محیط می‌ماند.
\item سپس، آنافیلاکسی، مسمومیت \ShE{CCB}، \ShU{آنتی‌سایکوتیک \ShE{II}}.
\item شوک نوروژنیک: قطع ایمپالس از سمپاتیک به عروق.
\item آسیب شدید تروماتیک مغز یا نخاع گردن.
\item چون فعالیت واگ هست: برادی‌کاردی یا \ShE{HR} طبیعی.
\end{itemize}

\ShPage{۳}
\subsection*{شوک دیسوسییتیو}
\begin{itemize}
\item خون‌رسانی بافتی \ShE{N/L} است ولی اکسیژن‌رسانی به سلول‌ها مختل است.
\item آنمی شدید، مت‌هموگلوبینمی، مسمومیت \ShE{CO}، سیانید (شایعین).
\end{itemize}
\subsection*{تشخیص شوک سپتیک در کودکان}
امتیاز حداقل ۲ در \ShE{PSS} (فونیکس سپس اسکور) از مجموع چهار ارگان: تنفسی، نورولوژیک، انعقادی، کاردیوواسکولار؛ که حداقل ۱ امتیاز آن مربوط به کاردیوواسکولار باشد.

در کودکان کم‌سن، شوک سپتیک ممکن است به صورت برون‌ده قلبی کاهش‌یافته و \ShM{SVR\uparrow} باشد.

\subsection*{شوک هایپوولمیک}
\begin{itemize}
\item برون‌ده قلبی \ShM{\downarrow}؛ پره‌لود \ShM{\downarrow}.
\item باید ۳۰\% حجم در گردش از دست برود تا کاهش فشار خون حاصل شود.
\item \ShE{Narrow Pulse Pressure}، \ShE{MSC}، \ShM{UOP\downarrow}.
\item درمان: مایع‌درمانی، وازوپرسور، تزریق خون.
\end{itemize}
\subsection*{شوک کاردیوژنیک}
\begin{itemize}
\item \ShM{CO\downarrow}، افترلود \ShM{\uparrow}.
\item انتهاهای سرد، \ShE{NPP}، دیسترس تنفسی (رال)، هپاتومگالی، \ShE{JVD}.
\item در \ShE{CXR}: کاردیومگالی؛ \ShU{کوآرکتاسیون}، تفاوت بین \ShE{BP} اندام فوقانی/تحتانی.
\item مایع‌درمانی محافظه‌کارانه، اینوتروپ، وازودیلاتور (برای کاهش افترلود).
\item رفع تاکی‌آریتمی‌ها؛ \ShE{PGE} برای زیر ۱ ماهه.
\end{itemize}

\ShPage{۴}
\subsection*{شوک توزیعی}
\begin{itemize}
\item \ShM{CO\downarrow}، \ShM{PL\downarrow}، \ShM{AL\downarrow}.
\item انتهاها گرم هستند؛ \ShE{WPP}، پالس \ShE{bounding}.
\item آنافیلاکسی: راش، تورم.
\item شوک نخاعی: ناتوانی در \ShM{HR\uparrow}، \ShM{BP\downarrow}.
\item نوروژنیک: \ShM{ICP\uparrow}.
\item مایع‌درمانی، وازودیلاتور.
\end{itemize}
\subsection*{شوک انسدادی}
انتهاها سرد، \ShM{CO\downarrow}، \ShM{AL\uparrow} بطن راست؛ \ShE{NPP}، وریدهای گردن برجسته، صداهای تنفسی غیرقرینه.

\subsection*{شوک دیسوسییتیو}
\begin{itemize}
\item انتهاها سرد.
\item جایگزینی مایع؛ خون باید تزریق باشد.
\item اکسیژن.
\end{itemize}
شوک در ابتدا ممکن است فقط به صورت یک تاکی‌کاردی \ShM{\pm} تاکی‌پنه بروز کند.

شوک هایپوولمیک: مخاطات خشک، اشک‌ها خشک، تورگور پوستی \ShM{\downarrow}، کاهش برون‌ده ادراری.

شوک کاردیوژنیک: تاکی‌پنه، انتهاهای سرد، \ShE{MSC}، نبض ضعیف.

\ShPage{۵}
در شوک توزیعی، هایپوتنشن دیاستولیک شایع‌تر است (نوروژنیک).

\subsection*{شوک سپتیک}
تب یا هایپوترمی وجود دارد.
\begin{itemize}
\item شوک سرد: شیرخوار و کودک کم‌سن؛ شبیه شوک هایپوولمیک و کاردیوژنیک.
\item شوک گرم: کودکان بزرگ‌تر و نوجوانان؛ شبیه شوک توزیعی.
\item ممکن است به هم تبدیل شوند.
\end{itemize}
هایپرفریتینمی (فریتین \ShM{>3000\;\mathrm{ng/ml}}) در شوک سپتیک:
\begin{itemize}
\item اندیکاسیون \ShE{PICU}.
\item ممکن است به علت \ShE{HLH} اولیه (معمولاً زیر ۲ سال) یا \ShE{HLH} ثانویه (شایع‌تر) باشد.
\end{itemize}
از \ShE{POCUS} برای ارزیابی بیمار و وضعیت مایعات و بررسی علل شوک استفاده می‌شود؛ دیلاتاسیون \ShE{RV}: \ShE{PE}.

\subsection*{آزمایشگاه}
حتماً \ShE{LDH}، حتماً \ShE{ABG}، الکترولیت‌ها، هایپوکلسمی، هایپوگلیسمی.

لاکتات شریان \ShM{>2\;\mathrm{mmol/L}}: پروگنوز بد.

\ShPage{۶}
\subsection*{ادامهٔ آزمایشگاه}
\begin{itemize}
\item \ShE{CBC}، \ShE{BUN}، \ShE{Cr}، \ShE{LFT}، انعقادی، فیبرینوژن، \ShE{D-dimer}، فریتین.
\item در عفونت ممکن است \ShM{CRP\uparrow} و پروکلسی‌تونین \ShM{\uparrow}.
\item نوتروفیل \ShM{\uparrow} و فرم‌های نابالغ \ShM{\uparrow} (باند، میلوسیت، پرومیلوسیت).
\item هم \ShE{ABG}، هم \ShE{VBG}: اختلاف \ShM{PCO_2} وریدی ـ شریانی \ShM{>5\;\mathrm{mmHg}}: \ShE{Early shock state}.
\item حداقل ۱ کشت خون.
\item فریتین در همهٔ بیماران؛ اگر بالای \ShU{۳۰۰؛ رقم نخست ممکن است ۵ باشد}: \ShE{MAS}، \ShE{HLH}؛ \ShE{TG}، \ShE{IL-2} رسپتور.
\end{itemize}
\subsection*{ملاحظات منیج شوک در کودکان}
\begin{itemize}
\item راه هوایی باز باشد؛ تنفس به خوبی انجام شود؛ گردش خون درست باشد.
\item اگر \ShM{SpO_2<92\text{--}94\%}: اکسیژن‌درمانی.
\item دسترسی وریدی بگیریم؛ ظرف دقایق اول.
\item ارزیابی مجدد بیمار در فواصل کم.
\item اگر ظرف ۵ دقیقه نتوانست \ShE{IV} بگیرد: \ShE{IO}.
\item اگر علی‌رغم مایع‌درمانی \ShE{40--60 ml/kg} بولوس: \ShE{CV-line}؛ \ShE{fluid-refractory shock}.
\end{itemize}

\ShPage{۷}
\begin{itemize}
\item اگر نیاز به وازواکتیو بود: \ShE{CV-line}.
\item اگر نیاز به دکستروز یا الکترولیت غلیظ: \ShE{CV-line}.
\end{itemize}
\subsection*{احیا با مایعات}
\begin{itemize}
\item دوزهای منقسم بولوس: \ShE{10--20 ml/kg}.
\item هدف: کاهش \ShE{HR}؛ افزایش \ShE{BP}؛ دمای انتهاها؛ پالس انتهاها؛ پرشدن مویرگی؛ ادرار.
\item افت بیش از ۱۰–۲۰\% \ShE{LDH}.
\item ممکن است طی ساعت اول \ShE{60 ml/kg} مایع نیاز داشته باشیم و ظرف چند ساعت اول تا \ShE{200 ml/kg}.
\item نحوهٔ مایع‌دادن باید \ShE{IV push} باشد.
\item هنگام مایع‌دادن بررسی از نظر \ShE{Overload}: شامل رال، ریتم گالوپ، هپاتومگالی.
\end{itemize}
اگر شواهد \ShE{Overload} رخ داد، حجم داخل عروق \ShE{OK} بود و همودینامیک علی‌رغم احیا بهبود پیدا نکرد: احیا با مایعات تعدیل؛ احیا با وازوکانستریکتور.

\ShPage{۸}
در شوک سپتیک، مایع کلوئیدی اندکی از کریستالوئیدی بهتر است: آلبومین ۴–۵\%.

محلول‌های کلوئید سنتتیک مثل ژلاتین و هیدروکسی‌اتیل‌استارچ ریسک \ShE{AKI} و کوآگولوپاتی دارند؛ توصیه نمی‌شوند.

چون کریستالوئید در دسترس و ارزان است، منیج اولیهٔ شوک با سرم کریستالوئیدی است.
\begin{itemize}
\item استثنا: شوک هموراژیک.
\item شوک دیسوسییتیو ناشی از آنمی: \ShE{Pack-cell}.
\end{itemize}
احیا حجم با \ShE{N/S (0.9\%)}: خطر اسیدوز متابولیک هایپرکلرمیک می‌شود.

در شوک هموراژیک، احیا با \ShE{P.C} بهتر از \ShE{N/S} است؛ نسبت \ShE{1:1:1} بین \ShE{P.C}، پلاسما و پلاکت.

در سایر انواع شوک اگر \ShE{Hb} زیر ۹ بود: \ShE{P.C}. \ShU{حروف قلم‌خوردهٔ کنار \ShE{Hb}: \ShE{P.C}.}

در این شرایط نیاز به \ShE{Plt} و پلاسما نیست مگر خونریزی.

\ShPage{۹}
\subsection*{داروهای وازواکتیو}
در کودکان با شوک راجعه علی‌رغم مایع‌درمانی:
\begin{itemize}
\item اینوتروپ‌ها مثل \ShE{EPN} (خط اول)، خط دوم دوپامین.
\item اندیکاسیون: \ShE{MI}، شوک کاردیوژنیک.
\item وازوپرسورها: \ShE{NEP}؛ زمانی که وازوپلژی وجود دارد؛ شوک توزیعی، نوروژنیک.
\item بهتر است ابتدا \ShE{CV-line} تزریق شود، ولی از \ShE{IV} و \ShE{IO} نیز ایمن.
\end{itemize}
\subsection*{الگوریتم برخورد با شوک سپتیک و سپس}
شک به \ShE{Sepsis}: ۳ ساعت فرصت بررسی داریم؛ در این مدت اگر تبدیل به شوک شد یا نارسایی ارگانی دیدیم، مثل شوک سپتیک برخورد. \ShU{انتهای سطر در لبهٔ چپ تصویر بریده است.}

\ShE{Septic Shock}: ظرف ۱ ساعت باید اقدامات زیر انجام شوند:
\begin{enumerate}
\item دسترسی وریدی مناسب.
\item \ShE{BC}.
\item شروع آنتی‌بیوتیک تجربی وسیع‌الطیف.
\item اندازه‌گیری سطح لاکتات.
\item مایع‌درمانی بولوس.
\item شروع داروی وازواکتیو در صورت نیاز.
\end{enumerate}

\ShPage{۱۰}
\subsection*{کودک با شوک سپتیک}
\subsubsection*{وجود \ShE{PICU}}
دیگر وجود یا نبود \ShM{BP\downarrow} مهم نیست.
\begin{itemize}
\item در نبود شواهد \ShE{Overload} مایع: بولوس \ShE{10--20 ml/kg}.
\item ارزیابی مجدد و دادن بولوس‌های مجدد تا رفع شوک یا ایجاد \ShE{Overload} مایع.
\item ارزیابی عملکرد قلبی.
\item اگر اختلال عملکرد میوکارد: \ShE{EPN}.
\item اگر شوک علی‌رغم \ShE{40--60 ml/kg} مایع رفع نشد یا شواهد \ShE{Overload} ظاهر شد: \ShE{EPN} یا \ShE{NEP}.
\end{itemize}
\subsubsection*{نبود \ShE{PICU}}
\textbf{عدم \ShM{BP\downarrow}:}
\begin{itemize}
\item اگر شواهد دهیدراتاسیون ندارد، بولوس ندهید (اسهال و غیره).
\item مایع \ShE{maintenance} بدهید.
\item مانیتورینگ همودینامیک.
\item در نظر گرفتن وازواکتیو ـ اینوتروپ.
\end{itemize}
\textbf{\ShM{BP\downarrow}:}
\begin{itemize}
\item در نبود شواهد \ShE{Overload}: بولوس \ShE{10--20 ml/kg}.
\item ارزیابی مجدد، بولوس مجدد.
\item ارزیابی عملکرد قلبی.
\item پس از \ShE{40 ml/kg}: \ShE{EPN} یا \ShE{NEP}.
\end{itemize}

\ShPage{۱۱}
پس از رفع شوک و بهبود پرفیوژن محیطی:
\begin{itemize}
\item دیگر مایع بولوس ندهید.
\item مایع \ShE{Maintenance} بدهید.
\item ارزیابی دوره‌ای از نظر بازگشت شوک.
\end{itemize}
\subsection*{معیارهای \ShM{BP\downarrow} و هایپوتنشن}
\begin{itemize}
\item سن زیر ۱ سال: \ShM{SBP<50\;\mathrm{mmHg}}.
\item سن ۱–۵ سال: \ShM{SBP<60\;\mathrm{mmHg}}.
\item سن بالای ۵ سال: \ShM{SBP<70\;\mathrm{mmHg}}.
\item یا انتهاهای سرد + پرشدن مویرگی بالای ۳ ثانیه + \ShE{weak/fast pulse}.
\end{itemize}
\subsection*{\ShE{EPN}}
\begin{itemize}
\item اینوتروپ و کرونوتروپ.
\item در دوزهای پایین وازودیلاتور است ولی در دوز بالای \ShE{0.1 mcg/kg/min} یک وازوکانستریکتور قوی.
\item دوز درمانی در شوک: \ShE{0.05--2 mcg/kg/min}.
\item در دوزهای بالا خون‌رسانی به کلیه را مختل می‌کند.
\item در دوز بالا ریسک آریتمی.
\end{itemize}
\subsection*{\ShE{NEP}}
\begin{itemize}
\item دوز درمانی در شوک مشابه \ShE{EPN}.
\item \ShE{BP} را ثانویه به افزایش مقاومت عروق محیطی بالا می‌برد.
\item \ShE{AF} بطن چپ را \ShM{\uparrow}.
\item اینوتروپ هم هست ولی اثر آن به دنبال \ShM{AF\uparrow} خنثی.
\end{itemize}

\ShPage{۱۲}
\subsection*{میلرینون}
\begin{itemize}
\item اینوتروپ مثبت؛ بهبود عملکرد دیاستولیک قلب؛ وازودیلاتاسیون عروق محیطی.
\item \ShE{0.25--1 mcg/kg/min}.
\end{itemize}
در آنمی شدید و شوک دیسوسییتیو به همه اکسیژن می‌دهیم:
\begin{latin}
\begin{center}
100\% \(O_2\) with non-rebreather mask
\end{center}
\end{latin}
\subsection*{اندیکاسیون انتوباسیون و \ShE{MV}}
\begin{itemize}
\item نارسایی تنفسی حاد هایپرکاپنیک یا هایپوکسیک.
\item باقی‌ماندن نارسایی \ShE{Circulatory} علی‌رغم درمان؛ برای کاهش \ShE{demand} متابولیک به‌وسیلهٔ \ShE{Sedation}.
\end{itemize}
ونتیلاسیون با فشار مثبت: کاهش بازگشت وریدی، که در بیماری که هنوز هایپوولمیک است می‌تواند به نارسایی قلب منجر شود.

بهترین دارو برای \ShE{Sedation}: کتامین؛ کمترین اثر بر همودینامیک. از اتومیدیت استفاده نکنید.

\subsection*{درمان اختصاصی برای شوک کاردیوژنیک}
علل: کاردیومیوپاتی، میوکاردیت، \ShE{MIS-C}، سایر علل.

\ShPage{۱۳}
\subsection*{ادامهٔ درمان شوک کاردیوژنیک}
\begin{itemize}
\item کاهش قدرت انقباضی قلب باید اصلاح شود.
\item درمان زودهنگام با اینوتروپ مثبت: \ShE{EPN}، دوپامین، دوبوتامین، میلرینون.
\item مایع‌درمانی خفیف برای بهبود پره‌لود.
\item ممکن است پرفیوژن پریفرال ضعیف، اسیدوز و \ShE{SVR} بالا علی‌رغم شروع اینوتروپ مثبت باقی بماند: شروع اینودیلاتور مثل میلرینون.
\item میلرینون عملکرد \ShE{S} را بهبود و \ShE{SVR} را کاهش می‌دهد ولی تأثیری بر \ShE{HR} ندارد؛ همچنین ریلکسیشن \ShE{D} را بهبود می‌دهد.
\item دوبوتامین و نیتروپروساید: خط‌های بعدی.
\end{itemize}
\subsection*{درمان اختصاصی شوک انسدادی}
\begin{itemize}
\item مداخله برای رفع انسداد به صورت مکانیکی.
\item در نوزادان و شیرخواران کوچک که به علت بیماری مادرزادی وابسته به داکت دچار شوک شده‌اند: انفوزیون \ShE{PG-E1}، \ShE{0.01--0.2 mcg/kg/min}.
\item وازودیلاتور؛ برون‌ده قلبی را حفظ می‌کند تا زمان انتقال به مرکز مجهز.
\end{itemize}
اگر شوک سپتیک روی داروی وازواکتیو رفع نشد:
\begin{itemize}
\item استرس دوز کورتیزون: هیدروکورتیزون.
\item بررسی علل آتیپیکال.
\item \ShE{ECMO}.
\end{itemize}

\ShPage{۱۴}
\subsection*{علل آتیپیک شوک}
هایپوتیروئیدی و سندرم کمپارتمان شکمی (هایپرتنشن داخل شکمی).

\subsection*{استرس دوز کورتیکواستروئید}
\ShE{50--100 mg/m\textsuperscript{2}/day}.
\begin{itemize}
\item بیمار شوک سپتیک که نارسایی مطلق آدرنال دارد.
\item نارسایی مطلق در زمینهٔ درمان گلوکوکورتیکوئید رخ می‌دهد.
\item \ShE{CAH} (هایپرپلازی مادرزادی).
\item شوک سپتیک با پورپورا.
\item درمان با استروئید قبلی.
\item هایپوپیتوئیتاریسم.
\item \ShE{Critical illness related corticosteroid insufficiency}.
\end{itemize}
\subsection*{اهداف درمان شوک}
\begin{itemize}
\item کاهش \ShE{HR}، کاهش \ShE{RR}، افزایش \ShE{BP}، افزایش ادرار به بالای \ShE{0.5 cc/kg/h}، نرمال‌شدن \ShE{MS}، گرم‌شدن انتهاها.
\item \ShE{LDH} کمتر از \ShE{2 mmol/L} طی ۴ ساعت عالی است.
\end{itemize}

\newcommand{\ShOriginal}[2]{%
  \clearpage\begin{center}\textbf{تصویر اصلی #1}\par
  \includegraphics[width=\textwidth,height=0.89\textheight,keepaspectratio]{images/pic#2.jpg}%
  \end{center}}
\ShOriginal{۱}{01}
\ShOriginal{۲}{02}
\ShOriginal{۳}{03}
\ShOriginal{۴}{04}
\ShOriginal{۵}{05}
\ShOriginal{۶}{06}
\ShOriginal{۷}{07}
\ShOriginal{۸}{08}
\ShOriginal{۹}{09}
\ShOriginal{۱۰}{10}
\ShOriginal{۱۱}{11}
\ShOriginal{۱۲}{12}
\ShOriginal{۱۳}{13}
\ShOriginal{۱۴}{14}
\end{document}
